\subsection{无处稠密集}

定义 1. 拓扑空间 X 的子集 A 称为无处稠密的，如果它的闭包没有内点。

等价地，A 在 X 中无处稠密当且仅当 A 的外部在 X 中稠密。

闭集 A 无处稠密当且仅当它没有内点，亦即当且仅当它与自己的边界相重合。任一子集 A 无处稠密当且仅当 A 的闭包无处稠密。无处稠密集的每个子集都无处稠密。

例。1) X 的空子集是无处稠密的。在豪斯多夫空间 X 中，由单独一点组成的集是无处稠密的，当且仅当该点在 X 中不是孤立点。稠密子集决不会是无处稠密的（除非 \(X = \varnothing\) ）。

\begin{enumerate}
\def\labelenumi{\arabic{enumi})}
\setcounter{enumi}{1}
\item
开集或闭集的边界总是无处稠密的。
\item
在空间 \(\mathbf{R}^n\) 中，每个维数 \(p < n\) 的线性子空间都是无处稠密集（第 VI 章，§ 1，第 4 号，命题 2）。
\end{enumerate}

注。任意子集的边界未必是无处稠密的：例如，若 A 与 \(\mathbb{C}\) A 都是稠密集，则 A 的边界是整个空间。

命题 1. 有限多个无处稠密集的并是无处稠密的。

只需证明两个无处稠密集 A、B 的并是无处稠密的，并且不失一般性可设 A 与 B 都是闭集。于是该命题等价于说两个稠密开集 CA、CB 的交是稠密的。现设 U 是非空开集，则 \(U \cap CA\) 是开集且非空，从而

\[
\left(\mathrm{U} \cap \mathbb {C A}\right) \cap \mathbb {C B} = \mathrm{U} \cap \left(\mathbb {C A} \cap \mathbb {C B}\right)
\]

是开集且非空。

设 Y 是拓扑空间 X 的子空间。Y 的子空间 A 称为相对于 Y 无处稠密的，如果 A 作为拓扑空间 Y 的子集来看是无处稠密的。

命题 2. 设 Y 是拓扑空间 X 的子空间，A 是 Y 的子集。若 A 相对于 Y 无处稠密，则 A 相对于 X 无处稠密。反之，若 Y 在 X 中开且 A 相对于 X 无处稠密，则 A 相对于 Y 无处稠密。

设 A 相对于 Y 无处稠密。若 A 在 X 中的闭包 \(\overline{A}\) 含有一个非空开集 U ，则 \(U \cap A\) 非空（由闭包的定义）；于是 \(U \cap Y\) 是相对于 Y 的非空开集，且含于 A 相对于 Y 的闭包 \(\overline{A} \cap Y\) 之中，这与假设相矛盾。

现设 Y 在 X 中开，且 \(A \subset Y\) 相对于 X 无处稠密。若 U 在 Y 中开且非空，则 U 在 X 中开，从而含有一个在 X 中（因而在 Y 中更是）开的非空集 V ，且 V 不与 A 相交；故 A 相对于 Y 无处稠密。

若 Y 在 X 中不开，命题 2 的第二部分显然不成立；例如考察 \(Y \neq \varnothing\) 在 X 中无处稠密且 A = Y 的情形。

